%% =================================================================
%% Wei, Mei, Zhao, Xiao, Sun, Zhang, Guo.
%% "Nondestructively identifying the mechanical behavior of soft
%%  tissues using surface deformation with an explicit inverse
%%  approach."
%% Appl. Math. Modelling 134 (2024) 126--147.
%%
%% Reconstruction of page 19 (printed p. 144) as study notes.
%% Generated by: reproduce-basic-paper-tex → reproduce-multiple-pages-basic-tex [17, 22]
%%
%% Structure: Table 5 (avg rel error, one-inclusion case) + Fig. 25
%% (Sample 1 with 5% noise) + 3 substantive Discussion paragraphs
%% (this is the most self-critical prose in the paper — the authors
%% explicitly acknowledge when the implicit method would outperform
%% theirs). No equations.
%% =================================================================

\documentclass[11pt]{article}

\usepackage[margin=1in]{geometry}
\usepackage{amsmath, amssymb}
\usepackage{mathrsfs}
\usepackage{bm}
\usepackage{xcolor}
\usepackage{soul}
\definecolor{hlyellow}{RGB}{255,255,170}
\definecolor{hlcyan}{RGB}{170,255,255}
\definecolor{hlgreen}{RGB}{170,255,170}
\definecolor{hlpink}{RGB}{255,170,170}
\definecolor{hlorange}{RGB}{255,221,170}
\usepackage{fancyhdr}
\usepackage{booktabs}
\usepackage{multirow}
\usepackage{array}

\pagestyle{fancy}
\fancyhf{}
\lhead{\small Y.~Mei et al.}
\rhead{\small Applied Mathematical Modelling 134 (2024) 126--147}
\cfoot{\small 144 $\to$ \arabic{page}}
\renewcommand{\headrulewidth}{0.4pt}

\setcounter{page}{1}

\begin{document}

% ---------- Table 5 ----------
\noindent\textbf{Table 5}\\[0.2em]
{\small Comparison of the average relative error in the shear
modulus distribution obtained by implicit inverse and explicit
inverse methods for the case with one inclusion embedded
(Figs.~24 and~25).}
\vspace{0.3em}

\begin{center}
\small
\begin{tabular}{@{}l ccc ccc@{}}
\toprule
Inverse method
  & \multicolumn{3}{c}{Explicit inverse method}
  & \multicolumn{3}{c}{Implicit inverse method} \\
\cmidrule(lr){2-4}\cmidrule(lr){5-7}
Number of surface displacement datasets
  & 3 & 6 & 9 & 3 & 6 & 9 \\
\midrule
1\,\% noise & 1.81\,\% & 1.50\,\% & 1.60\,\% & 4.28\,\% & 3.58\,\% & 3.41\,\% \\
5\,\% noise & 1.88\,\% & 1.86\,\% & 1.74\,\% & 6.49\,\% & 6.05\,\% & 4.26\,\% \\
\bottomrule
\end{tabular}
\end{center}

\vspace{1em}

% ---------- Figure 25 placeholder ----------
\noindent
\begin{center}
\fbox{\parbox{0.92\textwidth}{\centering\vspace{1em}
\textbf{[Figure 25 --- placeholder]}\\[0.8em]
Six-panel comparison of \textbf{Explicit Inverse Method} (top
row) vs.\ \textbf{Implicit Inverse Method} (bottom row) on
Sample 1 with \textbf{5\,\% noise} added, using 3, 6, and 9
displacement datasets.\\[0.4em]
\textbf{Explicit (top row)}:
(a) 3 fields, SM $\in [1.00, 5.83]$\,MPa --- 17\,\% overshoot;
(b) 6 fields, SM $\in [1.00, 6.09]$\,MPa --- \textbf{22\,\%
overshoot, WORSE than 3-field};
(c) 9 fields, SM $\in [1.00, 4.66]$\,MPa --- now 7\,\%
\emph{undershoot}. The max is non-monotone
($5.83 \to 6.09 \to 4.66$) and crosses the nominal 5\,MPa target
between 6 and 9 fields. This contradicts the intuitive ``more
data $\Rightarrow$ better'' story.\\[0.4em]
\textbf{Implicit (bottom row)}:
(d) 3 fields, SM $\in [1.00, 2.90]$\,MPa --- \textbf{42\,\%
undershoot};
(e) 6 fields, SM $\in [1.00, 3.21]$\,MPa;
(f) 9 fields, SM $\in [1.00, 3.40]$\,MPa --- still 32\,\%
undershoot.
The implicit method never approaches the 5\,MPa target under
5\,\% noise; even 9 fields leaves it severely below
nominal.\\[0.3em]
\textbf{Table 5 vs. this figure --- an important disconnect}:
Table 5 reports the explicit-9-fields-5\%-noise case as
\emph{1.74\,\% average relative error}, but panel (c) shows the
tumor-stiffness max at \emph{4.66 instead of 5.00 MPa}. Both
numbers are correct --- Table 5 averages over the whole 2D
domain (which is mostly 1\,MPa background where any method is
trivially correct), while the figure shows the peak inside the
tumor itself. A clinician looking at the ``$\approx 2\,\%$
error'' claim would expect a $\pm 0.1$\,MPa estimate of the
tumor; the figure shows a $-7\,\%$ bias in the quantity that
actually matters. The prose should have reported tumor-specific
error alongside the domain-average metric.
\vspace{1em}}}
\end{center}
\vspace{0.3em}
\begin{center}
\textbf{Fig. 25.} The reconstructed results with varying 3, 6
and 9 surface displacement datasets corresponding to Sample 1
(5\,\% noise is introduced into the measured data, Unit: MPa).
\end{center}

\vspace{1em}

% ---------- Discussion paragraphs ----------
\noindent
In the numerical cases presented in the paper, the total number
of predefined interfaces (for layered structures) or components
(for composites with embedded inclusions) in the initial guesses
is intentionally kept not much larger than the actual target.
This is because the topological information of the elastic
property distribution can often be determined beforehand based
on prior knowledge or observations. For example, in biomedical
imaging, characteristics such as the number, location, and rough
shape of tumors can be observed and identified. Similarly, for
typical layered biological tissues like skins or corneas, the
total number of layers has likely been well studied through in
vitro experiments or other research. Due to this available prior
knowledge, it is unnecessary to set a large number of
deformable components or interfaces in the initial guess for
the explicit inverse method. By constraining the initial guesses
to a reasonable and relevant number of interfaces or inclusions,
the optimization process can be more focused and efficient.
This can lead to faster convergence and more accurate
characterization of the elastic property distribution with a
reduced risk of convergence to a non-physical solution.

However, it should be noted that the results of this study do
not conclusively demonstrate that the explicit inverse method
outperforms the implicit inverse method in all cases. In
scenarios where the elastic property distribution lacks
significant heterogeneity or where no prior knowledge is
available, the explicit inverse method may indeed perform less
effectively compared to the implicit inverse method. This is
because the explicit inverse method relies on predefined
interfaces or inclusions, which may not accurately capture the
complexity of the underlying distribution. In such situations,
the implicit inverse method, with its greater flexibility and
ability to explore a wider range of possibilities, might yield
more accurate and robust results. The choice between the
explicit and implicit inverse methods, therefore, depends on
the specific characteristics of the problem and the availability
of prior information, ensuring the selection of the most
suitable approach for achieving meaningful and reliable results.

\pdfcomment[icon=Comment,color=cyan]{User input: did they use experimental data in the Wei paper?\\\\Claude comments: The answer is NO — they did NOT use experimental data. The paper "primarily focuses on numerical verification," meaning all results come from simulated/synthetic forward-problem data with known ground-truth shear modulus distributions, not from real tissue measurements. They explicitly cite "limited availability of experimental datasets" as the reason and note that ref [38] (Goenezen et al., 2021) validated the IMPLICIT method on experimental data, but their OWN explicit method has only been tested on synthetic numerical examples.}\sethlcolor{hlcyan}\hl{Due to the limited availability of experimental datasets, this
paper primarily focuses on numerical verification. It should be
noted that the implicit inverse method has demonstrated
successful applications in reconstructing the nonhomogeneous
shear modulus distribution of soft solids using experimental
datasets [38]. However, considering the superior performance of
the explicit inverse solver in numerical validation, we are
confident that it can achieve even better results when applied
to the experimental datasets.} In

% [page ends here — continues on page 20]

\end{document}
